{\white

{\Large Гравітаційні хвилі}

\begin{itemize}

\item Гравітація --- це загадкова сила, яка утримує великі об'єкти разом і зв'язує наші планети, зірки та галактики. Зв'язків між гравітацією та електромагнітним випромінюванням або іншими фізичними силами, такими як слабка сила, електромагнітна сила та сильна сила, не встановлено.

\item Теоретично гравітація викривляє простір і час. Фактично гравітація відповідає за згинання світла, що спостерігається на прикладі гравітаційного лінзування далеких зірок та галактик.

\item ``Гравітаційні хвилі'' з'являються як хвилі у просторі-часі, що рухаються зі швидкістю світла. Довго теоретизовані, гравітаційні хвилі були вперше виявлені у 2015 році і були спричинені злиттям двох чорних дір на відстані 1,4 мільярда світлових років.

\end{itemize}
}

